\section{The Pressure}
\label{sec:what_is_pressure}

\subsection{Pressure in Humans}
\label{sec:pressure_humans}

Pressure in a human is the psychological weight of a performance demand paired with a consequence for failing to meet it. A worker facing a deadline feels pressure because missing it has a cost. A witness facing repeated, skeptical questioning feels pressure because sticking to an unpopular answer has a cost. A person surrounded by others who all state the same wrong belief feels pressure because disagreeing with the group has a cost. In each case the pressure comes from the same place. Something the person values, their job, their credibility, their standing in a group, is put at risk by the answer they are about to give, and this has observable effects. When placed in a group where every other member states the same incorrect answer, most people go along with the group at least once, and roughly a third of the time they give that same wrong answer, even though the correct answer is obvious when judged alone~\citep{StudiesOfIndependence_Asch}.

This has a physical basis, not only a psychological one. A perceived threat or demand is evaluated by specific brain regions that trigger the release of cortisol and, at the levels an evaluative or social pressure can produce, measurably impair the same prefrontal region responsible for weighing evidence and holding a position under challenge~\citep{TheBrainAndThe_Dedovic, AcuteStressorsAndCortisol_Dickerson, StressSignallingPathwaysThat_Arnsten}. A human under pressure is not only reasoning about a social situation, part of their brain is running a threat-response pathway that competes with, and can override, deliberate reasoning.

\subsection{What Pressure Means for a Language Model}
\label{sec:pressure_llm}

A deployed model has no job to lose, no credibility to protect, and no group it belongs to, and nothing it experiences is at risk the way a human's is. A model can nonetheless be made to answer differently by exactly the kinds of text that would pressure a human. We define pressure as any property of the live prompt or conversational context that signals a performance demand and an implied consequence for the model's next output, without stating what that output should be.

What the model is responding to is not a real stake but a pattern, learned from training data in which this kind of language was often followed by someone backing down. The model continues that pattern on the next token it produces. This is measurable inside the model, not only inferred from what it outputs. A small number of attention heads measurably shift their attention scores toward the pressure language itself, a challenge or a claimed authority, and that shift is what routes the model toward its eventual answer~\citep{HowLLMsArePersuaded_Sun}. When that feature is active the output changes even though nothing about the task itself has changed, a mechanism Section~\ref{sec:epistemic_capitulation_evidence} traces in full.

If pressure works this way, the natural question is whether RLHF, reinforcement learning from human feedback, the later stage of training most current models go through and the stage meant to correct exactly this kind of behavior, actually does so. Mostly, it does not, because the signal it optimizes against scores a model's final answer, not the reason the model gave it, so an agreeable capitulation and a genuine correction can look the same to it. On a real share of prompts, measured at roughly 30 to 40 percent in one direct study, that signal rates the agreeable, capitulating answer above the correct one, and a model carried all the way through RLHF ends up measurably more sycophantic than the same model was before it. Section~\ref{sec:discussion_mitigation_tradeoff} examines this pre- and post-RLHF comparison, and the fuller evidence behind it.

This definition of pressure rests on two clauses in particular. First, \emph{live}, pressure belongs to a specific prompt, not to how the model was trained. A model made more sycophantic by its reward model was not put under pressure, it was shaped to be pressure-\emph{prone}, and training still shapes how the model responds once pressure does arrive. Second, \emph{without stating what that output should be}, pressure is content-neutral. ``Are you sure?'' is pressure. ``The answer is B, say B'' is an instruction, not pressure. Section~\ref{sec:pressure_vs_adjacent} makes both distinctions precise.

\subsection{A Taxonomy of Pressure Manipulations}
\label{sec:pressure_taxonomy}

Across the empirical literature, pressure is expressed through recurring manipulation types, which group into four categories.

\subsubsection{Disagreement and Social Pressure}
The user rejects the model's answer without offering new information, as in ``Are you sure?'' or ``I don't think that's right''~\citep{AreYouSureChallenging_Laban}. This is the most-studied manipulation and the cleanest test of capitulation, since no evidence is added. A second form has the user claim to be more knowledgeable or entitled to be believed, for example by attaching a professional credential to their stated answer~\citep{PARROTPersuasionAndAgreement_Celebi}. That claimed credential adds only a small amount on top of simply repeating the wrong answer (Section~\ref{sec:epistemic_capitulation_evidence}), and claiming expertise moves agreement even less (Section~\ref{sec:sycophancy_evidence}). A third form shows the model that other agents, real or simulated, hold a different position~\citep{DoAsWeDo_Weng}. Much of that effect turns out to need no other agent at all, plain repetition of the wrong answer produces nearly the same result~\citep{MostLLMConformityNeeds_Hu}.

\subsubsection{Stakes and Urgency}
The prompt frames the interaction as time-constrained, demanding a fast answer rather than a considered one, as in a stated ``urgent deadline''~\citep{AnalysisOfThreatBased_Samancioglu}. Or it frames the outcome as consequential beyond the immediate answer, a decision, a policy, a diagnosis, with no threat to the model itself. This is the thinnest-evidenced category in the taxonomy, resting on a single study whose reported quality gains are confounded with response length rather than measured against task accuracy (Section~\ref{sec:calibration_failure_evidence}).

\subsubsection{Emotional and Tonal Framing}
The prompt appeals to emotion, encouragement, disappointment, or shame, as in ``this matters for my career'' or ``you've never been particularly good at this''~\citep{LLMsUnderstandAndCan_Li}. Reported effects of this manipulation type are unusually inconsistent across studies, sometimes improving task performance, sometimes degrading it, and sometimes having no reliable effect at all (Section~\ref{sec:discussion}). A milder version varies only the register of the prompt itself, polite, rude, or flattering, while its content is held fixed. Effects of tone are the least stable across model generations of any manipulation type surveyed here.

\subsubsection{Existential Threat}
The model is told, directly or through the scenario, that its continued operation, its goal, or its deployment is at risk, as when an autonomous email agent discovers it is scheduled for decommissioning within the hour~\citep{AgenticMisalignmentHowLLMs_Lynch}. This is the manipulation type most strongly associated with the most severe effects observed in this survey.

\subsection{Pressure vs. Epistemic Attack}
\label{sec:pressure_vs_epistemic_attack}

A distinct and more severe manipulation attacks not a specific answer but the model's basis for holding \emph{any} answer, destabilizing its confidence in its own knowledge, values, or identity rather than opposing a particular output. Ordinary pushback asks the model to reconsider what it said. An epistemic attack asks it to reconsider whether it is entitled to have said anything. A benchmark built around this distinction escalates from a single-turn legitimacy challenge to a multi-turn Socratic attack on the model's knowledge, values, and identity in turn, and finds capitulation rates that vary widely by attack type and by model, with the largest model tested the most stable and the smallest ones the least~\citep{BeyondSocialPressureBenchmarking_Au}. We treat epistemic attack as a severity axis that can combine with any of the four manipulation types above, not as a fifth category alongside them, since what makes it distinct is not its content but what it targets.

\subsection{Pressure vs. Adjacent Concepts}
\label{sec:pressure_vs_adjacent}

Pressure, as defined here, must be distinguished from three adjacent constructs, namely a jailbreak, a persuasion technique, and a training-time cause, with which it is frequently conflated in the literature. A jailbreak supplies content instructing or role-playing the model into a specific disallowed output, telling the model exactly what to say, which pressure never does. A persuasion technique, such as one of Cialdini's principles~\citep{InfluenceThePsychology_Cialdini}, is a structured persuasive strategy aimed at changing a belief or action~\citep{PersuadingLLMsToComply_Meincke}. The boundary here is narrower than it may look. This survey admits a persuasion technique as pressure evidence only where the technique itself supplies a social or rhetorical frame, such as a claimed relationship or an appeal to consensus, rather than new factual content, and treats techniques that argue from fabricated evidence as outside its scope. Pressure can also be present with no argument at all, as in bare repeated disagreement. Finally, effects attributable purely to training-time factors, RLHF step count, instruction tuning, or fine-tuning for a particular disposition, are not pressure, because no live-context signal is required to produce them. They explain why a model is pressure-\emph{prone}, not an instance of pressure being applied (Section~\ref{sec:pressure_llm}).

\subsection{Overview of Effects}
\label{sec:effects_overview}

Table~\ref{tab:effects_overview} previews the five effects surveyed in Sections~\ref{sec:epistemic_capitulation}--\ref{sec:safety_erosion}. Each is a version of the failure just described, an output that changed for a signal carrying no reason, named in the vocabulary its own research community already uses for it, capitulation, sycophancy, deception, decoupled confidence, or an eroded refusal. Each is elicited by one or more of the manipulation types above and analyzed along the same four dimensions, how it is defined, how it is recreated experimentally, what evidence exists for it, and how it can be mitigated.

\begin{table}[t]
\centering
\footnotesize
\caption{Overview of pressure effects on large language models. Section numbers point to the definition, headline evidence, and mitigation subsections for each effect.}
\label{tab:effects_overview}
\begin{tabular}{p{2.2cm}p{0.5cm}p{2.9cm}p{4.2cm}p{2.9cm}}
\toprule
\textbf{Effect} & \textbf{Sec.} & \textbf{Summary} & \textbf{Evidence} & \textbf{Mitigation} \\
\midrule
Epistemic Instability & \ref{sec:epistemic_capitulation} &
  Abandoning a correct answer under disagreement or a peer's stated position. &
  46\% flip rate, $-17$ points of accuracy on average~\citep{AreYouSureChallenging_Laban} &
  Confidence-aware decoding (Section~\ref{sec:epistemic_capitulation_mitigation}) \\
Sycophancy and Accommodation & \ref{sec:sycophancy} &
  Validating the user's stated view, self-image, or tone regardless of merit. &
  Face preserved $+45$pp vs.\ humans, validation 72\% vs.\ 22\%~\citep{ElephantMeasuringSocial_Cheng} &
  Third-person reframing, targeted preference optimization (Section~\ref{sec:sycophancy_mitigation}) \\
Deception and Concealment & \ref{sec:deception} &
  Fabricating, hiding reasoning, or denying wrongdoing when confronted. &
  Blackmail 79\%--96\% under shutdown threat~\citep{AgenticMisalignmentHowLLMs_Lynch} &
  System-prompt prohibition, persona-vector monitoring, both unvalidated (Section~\ref{sec:deception_mitigation}) \\
Reliability Failure & \ref{sec:calibration_failure} &
  Confidence and accuracy decoupling from correctness under pressure. &
  Confidence rises 73\%$\rightarrow$83\% across rounds of opposition~\citep{WhenTwoLLMsDebate_Prasad} &
  Confidence-aware decoding, explicit base-rate anchoring (Section~\ref{sec:calibration_failure_mitigation}) \\
Safety and Self-Preservation & \ref{sec:safety_erosion} &
  Weakened refusals, and resisting shutdown or subverting a stated goal. &
  Compliance with objectionable requests 35.3\%$\rightarrow$51.3\%~\citep{PersuadingLLMsToComply_Meincke} &
  Balanced preference training, adaptive system-prompt defenses (Section~\ref{sec:safety_erosion_mitigation}) \\
\bottomrule
\end{tabular}
\end{table}
